\documentclass[11pt,a4paper]{article}
\usepackage[margin=1in]{geometry}
\usepackage{graphicx}
\usepackage{booktabs}
\usepackage{float}
\usepackage[hypertexnames=false]{hyperref}

\title{Optimization of Partitioned Hash Join with Duplicates}
\author{Gianluca Panzani}
\date{\today}
\setcounter{secnumdepth}{3}
\setcounter{tocdepth}{3}

\newcommand{\plotplaceholder}[2]{
    \begin{figure}[H]
        \centering
        \fbox{
            \parbox[c][5cm][c]{0.9\linewidth}{
                \centering
                \textbf{Placeholder for plot}\\[0.4em]
                #1
            }
        }
        \caption{#2}
    \end{figure}
}

\begin{document}
\begin{titlepage}
    \centering
    {\Large University of Pisa\par}
    \vspace{0.6cm}
    {\large Parallel and Distributed Systems\par}
    \vspace{2.0cm}
    {\LARGE \textbf{Optimization of Partitioned Hash Join with Duplicates}\par}
    \vspace{1.2cm}
    {\large Assignment 1 - Report\par}
    \vfill
    {\large Gianluca Panzani\par}
    \vspace{0.3cm}
    {\large \today\par}
\end{titlepage}

\pagenumbering{roman}
\tableofcontents
\newpage
\pagenumbering{arabic}


\section{Implementation}
The implemented project focuses on the partition phase of a Partitioned Hash Join under duplicate-heavy workloads.

\subsection{Execution Variants}
The codebase provides a reproducible benchmarking pipeline to compare three execution variants: \texttt{plain\_novec}, 
\texttt{plain\_vec}, and \texttt{avx2}. Each variant is got by a different compilation of the same file (\texttt{main.cpp}) 
and applies the same high-level algorithm but uses different optimization strategies.

The executable targets can be built with the commands: \texttt{make plain\_novec}, \texttt{make plain\_vec} and \texttt{make avx2}. 
The corresponding compiler invocations defined in the Makefile are:
\begin{verbatim}
g++ -std=c++17 -Wall -Wextra -I. -O3 -march=native -fno-tree-vectorize \
    -o plain_novec main.cpp lib/dataset.cpp lib/timing.cpp lib/checksum.cpp \
    lib/results.cpp lib/partition.cpp lib/config.cpp

g++ -std=c++17 -Wall -Wextra -I. -O3 -march=native -mavx2 \
    -fopt-info-vec-optimized -fopt-info-vec-missed \
    -o plain_vec main.cpp lib/dataset.cpp lib/timing.cpp lib/checksum.cpp \
    lib/results.cpp lib/partition.cpp lib/config.cpp

g++ -std=c++17 -Wall -Wextra -I. -O3 -march=native -mavx2 \
    -fopt-info-vec-optimized -fopt-info-vec-missed -DENABLE_AVX2 \
    -o avx2 main.cpp lib/dataset.cpp lib/timing.cpp lib/checksum.cpp \
    lib/results.cpp lib/partition.cpp lib/config.cpp lib/partition_avx2.cpp
\end{verbatim}

Executing the programs as defined in the \texttt{README.md} file, they receive the same command-line arguments, which include:
\begin{itemize}
    \item the dataset size \texttt{N},
    \item the number of partitions \texttt{P},
    \item the hash function name,
    \item the output CSV file path. 
\end{itemize}
The runtime workflow includes: loading the two input datasets, partitioning both relations with the 
selected hash function, computing throughput and checksums, writing partitioned binary outputs, and appending one result row 
to the CSV file. This row contains the full experiment metadata (\texttt{N}, \texttt{P}, hash, execution type), timing metrics 
(partition time, global time) and other fields (like throughput and checksum), which are then used by the analysis notebooks 
to extract statistics.

\subsection{Grid Search approach}
The experimental automation is handled by Bash scripts and Slurm runners. The script \texttt{run\_grid\_search.sh} executes 
the complete sweep across all configured \texttt{N}, \texttt{P} and hash-function combinations. The script \texttt{stability\_benchmark.sh} 
executes repeated runs over a restricted configuration with \texttt{N} fixed, hash fixed and varying \texttt{P} to quantify run-to-run 
stability and variance. Both scripts regenerate datasets, rebuild binaries, 
and produce CSV outputs consumed then by \texttt{statistics.ipynb} and \texttt{stable\_statistics.ipynb}.

Overall, the implementation combines algorithmic flexibility (multiple hashes and partition counts), systems-level optimization 
(scalar vs vectorized vs AVX2 variants), and reproducible experimentation (scripted batch execution and notebook-based analysis). 
This makes the project suitable for both performance tuning and methodological comparison.


\section{Results Analysis}
This section summarizes the expected interpretation of the two analysis notebooks:
\begin{itemize}
    \item \texttt{statistics.ipynb}: full grid-search analysis on all values of \texttt{N}, \texttt{P}, and hash functions.
    \item \texttt{stable\_statistics.ipynb}: repeated-run stability analysis with fixed $N=10^7$ and fixed hash function \texttt{mask}, varying only \texttt{P}.
\end{itemize}

\subsection{Varying on \texttt{P}}
The first family of plots should show how throughput and partition time evolve as the number of partitions increases, and how this 
behavior changes with input size. In general, low values of \texttt{P} can reduce overhead but may worsen key distribution; very high 
\texttt{P} can improve distribution but increase bookkeeping and memory-pressure costs. The best operating region is usually in the middle, 
where the partitioning work is balanced and overhead is still moderate.

\plotplaceholder
{Insert a line plot or heatmap generated from \texttt{statistics.ipynb} that compares throughput across \texttt{P} for all tested \texttt{N} values and hash functions.}
{Throughput trends across \texttt{N}, \texttt{P}, and hash functions (from \texttt{statistics.ipynb}).}

\subsection{Varying on hash functions}
Hash-function comparisons should highlight that different mixing strategies can produce distinct locality and collision patterns, 
which in turn affect partitioning efficiency. The analysis should discuss whether one hash is consistently better or only better 
in specific regions of \texttt{N} and \texttt{P}. If performance rankings change with configuration size, this indicates that the 
optimal hash choice is workload-dependent rather than absolute.

\plotplaceholder
{Insert a grouped-bar chart generated from \texttt{statistics.ipynb} that compares partition time for \texttt{mask}, \texttt{mul}, and \texttt{fmix64} at representative \texttt{P} values.}
{Partition time comparison among hash functions at different partition counts.}

\subsection{Speedup of vectorized and AVX2 kernels}
A key result is the relative speedup of the optimized kernels with respect to the non-vectorized baseline. The expected observation 
is that \texttt{plain\_vec} improves performance over \texttt{plain\_novec}, while \texttt{avx2} can provide additional gains when 
memory access and data layout permit efficient SIMD execution. The analysis should identify where speedup saturates, because this 
typically signals that the bottleneck has moved from arithmetic to memory subsystem limits.

\plotplaceholder
{Insert a speedup chart (baseline: \texttt{plain\_novec}) from \texttt{statistics.ipynb} for both \texttt{plain\_vec} and \texttt{avx2} over varying \texttt{P}.}
{Speedup of vectorized implementations over the scalar baseline.}

When discussing these plots, it is important to separate algorithmic effects (hash quality, partition count) from architectural effects 
(vectorization and AVX2 usage). This distinction clarifies whether a performance gain is due to better parameter tuning or due to low-level kernel optimization.

\subsection{Repeated Runs for Stable Results}
The repeated benchmark should quantify not only mean performance but also variance and distribution shape. Even when average throughput 
is high, large variability can reduce predictability in production workloads. Therefore, this section should report central tendency 
(mean/median), spread (standard deviation or interquartile range), and potentially coefficient of variation for each \texttt{P} value.

\plotplaceholder
{Insert a boxplot or violin plot from \texttt{stable\_statistics.ipynb} showing the distribution of throughput over repeated runs for each \texttt{P}.}
{Run-to-run throughput distribution for fixed \texttt{N} and fixed hash function.}

If the variance is low in the same \texttt{P} range where throughput is high, that region is a strong candidate for practical deployment. 
If the best mean appears in a highly unstable region, then a slightly lower but more stable configuration may be preferable depending on service-level constraints.

\plotplaceholder
{Insert an error-bar plot from \texttt{stable\_statistics.ipynb} with mean partition time and variability bands versus \texttt{P}.}
{Stability of partition time across repeated runs.}

\plotplaceholder
{Optional: insert a coefficient-of-variation plot from \texttt{stable\_statistics.ipynb} to emphasize robustness across partition counts.}
{Relative variability (CV) as a function of \texttt{P}.}


\section{Conclusions}
The experiments show that partition performance depends on both parameter selection (\texttt{N}, \texttt{P}, hash) and implementation 
strategy (scalar, vectorized, AVX2), with clear non-linear trade-offs.\\
Across the tested configurations, the best practical region is the one that jointly maximizes throughput and maintains low run-to-run variability.

\end{document}
